\section{Cơ sở lý thuyết}

\subsection{Phát hiện video độc hại}

Phát hiện video độc hại là một bài toán phân loại nội dung, trong đó hệ thống ước lượng liệu một video có chứa hoặc cổ xúy nội dung không an toàn, lạm dụng, gây hiểu lầm, bóc lột hoặc vi phạm chính sách hay không. Trong bối cảnh kiểm duyệt thực tế, mô hình loại này thường không được kỳ vọng tự đưa ra quyết định thực thi cuối cùng. Thay vào đó, mô hình có thể hỗ trợ phân loại ưu tiên bằng cách xác định những nội dung cần được rà soát kỹ hơn.

Đề tài này mô hình hóa nhiệm vụ cốt lõi dưới dạng phân loại nhị phân:

\begin{equation}
    y =
    \begin{cases}
        1, & \text{nếu video độc hại},\\
        0, & \text{nếu video không độc hại}.
    \end{cases}
\end{equation}

Cách phát biểu nhị phân được chọn làm sản phẩm chính vì đơn giản hơn, dễ đánh giá một cách đáng tin cậy hơn và phù hợp hơn để xây dựng một pipeline dữ liệu đầu cuối ổn định. Một cách phát biểu đa nhãn chi tiết hơn có thể được áp dụng khi cần gán video vào các loại tác hại cụ thể. Trong đề tài này, sáu nhãn loại \texttt{IH}, \texttt{HH}, \texttt{CB}, \texttt{ADD}, \texttt{SXL} và \texttt{PH} được xem là phần mở rộng giai đoạn hai, không thuộc hệ thống cốt lõi đã hoàn thiện.

\subsection{Thiết kế đặc trưng ưu tiên văn bản}

Video có thể chứa nhiều nguồn thông tin khác nhau, bao gồm khung hình, âm thanh, siêu dữ liệu, bình luận, tín hiệu tương tác và phần mô tả bằng văn bản. Đề tài này bắt đầu bằng một thiết kế ưu tiên văn bản (text-first). Đầu vào của mô hình được giới hạn trong ba trường:

\begin{itemize}
    \item \textbf{Tiêu đề}: phần tóm tắt ngắn được viết cho người xem.
    \item \textbf{Mô tả}: phần văn bản ngữ cảnh dài hơn đi kèm video.
    \item \textbf{Bản ghi lời nói}: nội dung nói được chuyển thành văn bản khi có sẵn.
\end{itemize}

Văn bản đầu vào hợp nhất cho mô hình có thể được biểu diễn như sau:

\begin{verbatim}
model_text = normalize(title)
           + normalize(description)
           + normalize(transcript)
\end{verbatim}

Các trường rỗng được bỏ qua và khoảng trắng được chuẩn hóa trước khi đưa vào mô hình. Thiết kế này giúp hệ thống có khả năng tái lập và giảm phụ thuộc vào các công cụ bên ngoài để trích xuất ảnh, xử lý âm thanh hoặc thu thập dữ liệu theo đặc thù nền tảng. Thiết kế này cũng giảm rủi ro rò rỉ dữ liệu bằng cách loại trừ các đặc trưng như danh tính kênh, số lượt xem và thời lượng. Danh tính kênh có thể khiến mô hình học các mẫu ở cấp người tạo nội dung thay vì tín hiệu ở cấp nội dung, trong khi các trường về độ phổ biến và thời lượng không có mặt nhất quán trên tất cả các nguồn có nhãn.

\subsection{Các tầng dữ liệu Bronze, Silver và Gold}

Học máy đáng tin cậy phụ thuộc vào kỹ thuật dữ liệu đáng tin cậy. Đề tài sử dụng mô hình tổ chức dữ liệu ba tầng: bronze, silver và gold. Mỗi tầng giữ một vai trò riêng trong pipeline.

\begin{table}[H]
\centering
\renewcommand{\arraystretch}{1.35}
\begin{tabularx}{\textwidth}{|>{\raggedright\arraybackslash}p{2.2cm}|>{\raggedright\arraybackslash}X|}
\hline
\textbf{Tầng} & \textbf{Mục đích} \\
\hline
Bronze & Bảo toàn các bản ghi nguồn thô và siêu dữ liệu nhập dữ liệu gần nhất có thể với tệp gốc. \\
\hline
Silver & Chuẩn hóa các nguồn không đồng nhất thành một lược đồ ở cấp dòng với tên trường, nhãn, trường văn bản và cờ chất lượng nhất quán. \\
\hline
Gold & Xây dựng các bảng sẵn sàng cho mô hình, trường văn bản chuẩn hóa, tập train/validation/test, tóm tắt đồng thuận và kết quả kiểm toán. \\
\hline
\end{tabularx}
\caption{Vai trò khái niệm của các tầng dữ liệu bronze, silver và gold.}
\label{tab:theory-data-layers}
\end{table}

Tầng bronze hỗ trợ khả năng truy vết vì bảo toàn nội dung đã được nhập. Tầng silver hỗ trợ tính nhất quán vì mọi nguồn đều được chuyển đổi về cùng một lược đồ. Tầng gold hỗ trợ mô hình hóa vì tạo ra một bản ghi chuẩn hóa cho mỗi mã định danh video và chuẩn bị đầu vào sạch cho huấn luyện cũng như đánh giá.

Thiết kế phân tầng đặc biệt hữu ích khi các tệp nguồn không đồng nhất. Trong đề tài này, dữ liệu bao gồm bảng tính, tệp CSV, nhiều nguồn chú thích, tệp trùng lặp, mã định danh video bị thiếu, tên cột không nhất quán và mức độ sẵn có của văn bản khác nhau. Pipeline phân tầng giúp các vấn đề này trở nên quan sát được thay vì bị che giấu bên trong mã huấn luyện mô hình.

\subsection{Giám sát yếu}

Giám sát yếu (weak supervision) là một chiến lược học sử dụng các nguồn nhãn chưa hoàn hảo để tăng lượng dữ liệu huấn luyện. Các nhãn này có thể rẻ hơn, nhiễu hơn hoặc ít có thẩm quyền hơn nhãn chuyên gia. Trong đề tài này, nhãn của chuyên gia lĩnh vực được xem là ground truth chính, trong khi nhãn của GPT-4-Turbo và crowdworker được xem là các nguồn yếu.

Ý tưởng cốt lõi không phải là tin tưởng mọi nhãn yếu như nhau. Thay vào đó, hệ thống ước lượng độ tin cậy của từng nguồn yếu bằng cách so sánh nguồn đó với nhãn chuyên gia lĩnh vực trên phần giao nhau thuộc tập huấn luyện có nhãn chuyên gia:

\begin{equation}
    r_s = \frac{\#\text{số lần nguồn }s\text{ đồng thuận với chuyên gia lĩnh vực}}{\#\text{số mẫu huấn luyện giao nhau}}.
\end{equation}

Các độ tin cậy sau đó được chuẩn hóa thành trọng số:

\begin{equation}
    w_{\text{GPT}} =
    \frac{r_{\text{GPT}}}{r_{\text{GPT}} + r_{\text{Crowd}}},
    \qquad
    w_{\text{Crowd}} =
    \frac{r_{\text{Crowd}}}{r_{\text{GPT}} + r_{\text{Crowd}}}.
\end{equation}

Đối với một video có các nhãn yếu khả dụng, xác suất độc hại yếu được tính như một trung bình có trọng số:

\begin{equation}
    p_{\text{weak}} =
    \frac{\sum_{s \in S} w_s y_s}{\sum_{s \in S} w_s},
\end{equation}

trong đó \(S\) là tập các nguồn yếu khả dụng và \(y_s \in \{0,1\}\) là nhãn nhị phân từ nguồn \(s\). Sau đó, hệ thống áp dụng một quy tắc chấp nhận thận trọng:

\begin{verbatim}
if weak_p >= 0.75:
    accept as harmful
elif weak_p <= 0.25:
    accept as harmless
else:
    abstain
\end{verbatim}

Các nhãn yếu được chấp nhận sẽ nhận trọng số mẫu dựa trên độ tự tin:

\begin{equation}
    \text{trọng số mẫu} = 0.5 + |p_{\text{weak}} - 0.5|.
\end{equation}

Quy tắc này gán trọng số cao hơn cho các nhãn yếu nằm xa vùng bất định và trọng số thấp hơn cho các nhãn gần biên quyết định.

\subsection{Gán nhãn giả}

Gán nhãn giả (pseudo-labeling) là một phương pháp tự huấn luyện, trong đó mô hình đã huấn luyện được dùng để gán nhãn cho các mẫu chưa có nhãn. Những nhãn được sinh tự động này chỉ được thêm vào tập huấn luyện nếu thỏa mãn các yêu cầu về độ tin cậy. Mục tiêu là khai thác tập chưa gán nhãn mà không cần chú thích thủ công cho từng mẫu.

Trong đề tài này, gán nhãn giả được thiết kế một cách thận trọng. Trước hết, mô hình giám sát yếu tốt nhất được dùng để chấm điểm các video chưa gán nhãn. Một nhãn giả chỉ được chấp nhận nếu mẫu có đủ văn bản và xác suất độc hại đã hiệu chỉnh nằm rất gần 0 hoặc 1:

\begin{verbatim}
if token_count >= 30 and (p_harmful >= 0.95 or p_harmful <= 0.05):
    accept pseudo-label
else:
    discard
\end{verbatim}

Tất cả nhãn giả được chấp nhận nhận trọng số mẫu cố định bằng 0.5. Mô hình được huấn luyện lại một lần duy nhất và quá trình tự huấn luyện dừng lại sau vòng đó. Cách làm này ngăn các vòng phản hồi không kiểm soát, trong đó lỗi của mô hình có thể bị củng cố lặp đi lặp lại.

Quy tắc cuối cùng cho gán nhãn giả dựa trên tập validation. Mô hình có nhãn giả chỉ được giữ lại nếu không làm giảm macro-F1 trên validation hoặc recall của lớp độc hại so với mô hình giám sát yếu. Quy tắc này quan trọng vì gán nhãn giả có thể làm tăng kích thước tập huấn luyện nhưng vẫn làm giảm chất lượng mô hình nếu các nhãn được chấp nhận bị lệch hoặc sai.

\begin{figure}[H]
\centering
\resizebox{0.98\textwidth}{!}{
\begin{tikzpicture}[
    box/.style={draw, rounded corners, align=center, minimum width=3.3cm, minimum height=1.05cm, font=\small},
    source/.style={box, fill=blue!10},
    model/.style={box, fill=green!10},
    decision/.style={box, fill=orange!15},
    arrow/.style={->, thick}
]
\node[source] at (0,2.4) (expert) {Nhãn chuyên gia\\lĩnh vực};
\node[source] at (0,0.8) (weak) {Nguồn yếu\\GPT và crowd};
\node[source] at (0,-0.8) (unlabeled) {Tập văn bản\\chưa gán nhãn};

\node[model] at (4.2,2.4) (baseline) {Mô hình nền\\chỉ chuyên gia};
\node[model] at (4.2,0.8) (weakmodel) {Chuyên gia +\\giám sát yếu};
\node[model] at (4.2,-0.8) (pseudomodel) {Chuyên gia + yếu +\\gán nhãn giả};
\node[decision] at (8.5,0.8) (selection) {Chọn mô hình\\theo validation};

\draw[arrow] (expert) -- (baseline);
\draw[arrow] (expert) -- (weakmodel);
\draw[arrow] (weak) -- node[above, font=\scriptsize] {nhãn có trọng số} (weakmodel);
\draw[arrow] (weakmodel) -- node[above, font=\scriptsize] {chấm điểm văn bản chưa gán nhãn} (pseudomodel);
\draw[arrow] (unlabeled) -- node[below, font=\scriptsize] {nhãn độ tin cậy cao} (pseudomodel);
\draw[arrow] (baseline) -- (selection);
\draw[arrow] (weakmodel) -- (selection);
\draw[arrow] (pseudomodel) -- (selection);

\node[align=center, font=\scriptsize] at (8.5,-1.0) {Chọn theo macro-F1\\và recall trên validation};
\end{tikzpicture}
}
\caption{Quan hệ khái niệm giữa học có giám sát, giám sát yếu, gán nhãn giả và lựa chọn mô hình dựa trên validation.}
\label{fig:theory-learning-strategy}
\end{figure}

\subsection{Biểu diễn văn bản bằng TF-IDF}

Các mô hình học máy cần đầu vào dạng số, vì vậy văn bản phải được chuyển thành vector đặc trưng. Đề tài sử dụng term frequency-inverse document frequency, thường được gọi là TF-IDF. TF-IDF gán trọng số cao hơn cho những từ xuất hiện thường xuyên trong một văn bản nhưng không quá phổ biến trên toàn bộ tập văn bản.

Với một thuật ngữ \(t\) trong văn bản \(d\), ý tưởng có thể được viết như sau:

\begin{equation}
    \text{TF-IDF}(t,d) =
    \text{TF}(t,d) \times \text{IDF}(t).
\end{equation}

Term frequency đo mức độ thường xuyên của một thuật ngữ trong một văn bản. Inverse document frequency làm giảm trọng số của các thuật ngữ xuất hiện trong nhiều văn bản và do đó kém phân biệt hơn. Đối với phát hiện video độc hại, điều này hữu ích vì các từ chung chung nên có ảnh hưởng thấp hơn những cụm từ hoặc thuật ngữ đặc thù cho ngữ cảnh độc hại hoặc không độc hại.

Hệ thống sử dụng cả unigram và bigram. Unigram biểu diễn từng từ đơn lẻ, trong khi bigram biểu diễn các cụm hai từ ngắn. Cách làm này giúp mô hình nắm bắt các cụm từ có ý nghĩa phụ thuộc vào tổ hợp từ thay vì từng từ rời rạc.

\subsection{Hồi quy logistic cho phân loại nhị phân}

Hồi quy logistic là một mô hình phân loại tuyến tính ước lượng xác suất của một kết quả nhị phân. Với vector đặc trưng \(x\), mô hình tính:

\begin{equation}
    P(y=1 \mid x) = \sigma(w^\top x + b),
\end{equation}

trong đó \(w\) là vector trọng số được học, \(b\) là hệ số chệch và \(\sigma\) là hàm sigmoid:

\begin{equation}
    \sigma(z) = \frac{1}{1 + e^{-z}}.
\end{equation}

Hồi quy logistic là một mô hình nền mạnh cho đặc trưng văn bản thưa vì hiệu quả, ổn định và dễ diễn giải hơn so với các mô hình neural nặng hơn. Mô hình này cũng phù hợp với quy trình cục bộ trên một máy đơn. Trong đề tài này, hiệu chỉnh xác suất (probability calibration) được áp dụng để các xác suất dự đoán phù hợp hơn cho việc đặt ngưỡng, ánh xạ mức rủi ro và chấp nhận nhãn giả.

\subsection{Các thước đo đánh giá}

Đề tài đánh giá phân loại nhị phân bằng nhiều thước đo thay vì chỉ dựa vào accuracy. Điều này quan trọng vì các bộ dữ liệu video độc hại có thể mất cân bằng, và một mô hình có thể trông chính xác về tổng thể nhưng lại hoạt động kém trên một lớp.

\begin{table}[H]
\centering
\renewcommand{\arraystretch}{1.3}
\begin{tabularx}{\textwidth}{|>{\raggedright\arraybackslash}p{3cm}|>{\raggedright\arraybackslash}X|}
\hline
\textbf{Thước đo} & \textbf{Ý nghĩa} \\
\hline
Precision & Trong số các video được dự đoán là độc hại, tỷ lệ video thực sự độc hại. \\
\hline
Recall & Trong số các video thực sự độc hại, tỷ lệ video được mô hình phát hiện. \\
\hline
Macro-F1 & Trung bình F1 trên các lớp, trong đó mỗi lớp có trọng số như nhau. \\
\hline
Weighted F1 & F1 được trung bình theo các lớp và có xét đến tần suất của từng lớp. \\
\hline
AUROC & Đo chất lượng xếp hạng của mô hình trên các ngưỡng phân loại khác nhau. \\
\hline
PR-AUC & Đo đánh đổi precision-recall, đặc biệt hữu ích khi dữ liệu mất cân bằng. \\
\hline
Ma trận nhầm lẫn & Đếm true negative, false positive, false negative và true positive. \\
\hline
\end{tabularx}
\caption{Các thước đo chính dùng để đánh giá phân loại nhị phân trong đề tài.}
\label{tab:theory-evaluation-metrics}
\end{table}

Macro-F1 và recall của lớp độc hại được nhấn mạnh trong đề tài. Macro-F1 ngăn lớp đa số chi phối kết quả đánh giá, trong khi recall của lớp độc hại cho biết mô hình phát hiện được bao nhiêu video độc hại. Precision cũng quan trọng vì false positive có thể khiến video không độc hại bị gắn cờ sai, nhưng đề tài ưu tiên recall và macro-F1 cho mục tiêu phân loại ưu tiên kiểm duyệt.

\subsection{Độ đồng thuận giữa các nguồn chú thích}

Vì bộ dữ liệu có nhãn từ chuyên gia lĩnh vực, GPT-4-Turbo và crowdworker, việc đo mức độ đồng thuận giữa các nguồn là cần thiết. Hai thước đo đồng thuận được sử dụng là tỷ lệ đồng thuận phần trăm và hệ số kappa của Cohen.

Tỷ lệ đồng thuận phần trăm là tỷ lệ trực tiếp các mẫu giao nhau mà hai nguồn gán cùng một nhãn:

\begin{equation}
    \text{Tỷ lệ đồng thuận} =
    \frac{\#\text{số nhãn khớp}}{\#\text{số mẫu giao nhau}}.
\end{equation}

Tuy nhiên, tỷ lệ đồng thuận phần trăm không tính đến mức đồng thuận có thể xảy ra do ngẫu nhiên. Hệ số kappa của Cohen hiệu chỉnh phần đồng thuận ngẫu nhiên:

\begin{equation}
    \kappa =
    \frac{p_o - p_e}{1 - p_e},
\end{equation}

trong đó \(p_o\) là mức đồng thuận quan sát được và \(p_e\) là mức đồng thuận kỳ vọng do ngẫu nhiên. Giá trị kappa thấp có thể cho thấy hai nguồn có tính nhất quán yếu sau khi hiệu chỉnh theo ngẫu nhiên, ngay cả khi tỷ lệ đồng thuận thô có vẻ ở mức trung bình.

Trong đề tài này, phân tích đồng thuận không chỉ có ý nghĩa mô tả. Kết quả này trực tiếp hỗ trợ thiết kế giám sát yếu bằng cách cho thấy nhãn của GPT và crowdworker nên được xem là nguồn nhiễu, không phải ground truth có độ tin cậy ngang bằng với chuyên gia.
